\documentclass[11pt,a4paper,landscape]{article}
\usepackage[T1]{fontenc}\usepackage[utf8]{inputenc}\usepackage[french]{babel}
\usepackage{lmodern,amsmath,amssymb,graphicx,booktabs,xcolor,hyperref,fancyhdr}
\usepackage[margin=16mm]{geometry}
\hypersetup{colorlinks=true,linkcolor=blue,urlcolor=blue,pdftitle={Plaque trouée : supports fins, modèle quadratique et continuation}}
\graphicspath{{../results/hudge_adaptive_support_20261115/figures/}}
\pagestyle{fancy}\fancyhf{}\fancyhead[L]{WGPU-Matching — Supports fins et courbure locale}
\fancyhead[R]{15 novembre 2026}\fancyfoot[R]{\thepage}\setlength{\headheight}{14pt}
\setlength{\parindent}{0pt}\setlength{\parskip}{6pt}
\newcommand{\fig}[2]{\begin{center}\includegraphics[width=\linewidth,height=#2,keepaspectratio]{#1}\end{center}}
\begin{document}
{\LARGE\bfseries Plaque trouée : faut-il réduire les fenêtres\par
ou ajouter une courbure à la transformation ?\par}
\vspace{4mm}
\textbf{Résultat : les deux pistes ont un intérêt distinct.} Le modèle quadratique
réduit nettement le résidu photométrique sur les mêmes supports de 32 pixels. La grille
affine fine apporte beaucoup plus de points réellement calculés. Mais une part des
échecs précédents était également due à la limite de 30 itérations, pas à une impossibilité
de représenter le déplacement par une affine locale.

Les 191 cibles sont comparées séparément à la référence. Le point de départ commun
est la reprise par fonctions de base radiale (RBF), contrainte par les gradients,
au seuil de cohérence 0,50 : 6\,549 centres retenus et \textbf{327 centres restants} sur 6\,876.
Les essais portent exclusivement sur les supports de ces centres restants.

\begin{center}\begin{tabular}{lrrr}
\toprule
Modèle & Centres récupérés, un appel & Centres récupérés, continuation & Points retenus, continuation\\\midrule
Affine $32\times32$, pas 32 &295/327&327/327&327/327\\
Quadratique $32\times32$, courbure initiale nulle &181/327&327/327&327/327\\
Quadratique $32\times32$, courbure initiale RBF &161/327&327/327&327/327\\
Affine $16\times16$, pas 8 &294/327&327/327&2\,843/2\,943\\\bottomrule
\end{tabular}\end{center}

Un appel comporte au maximum 30 itérations. La continuation relance uniquement les
sorties au statut non convergé, avec leurs propres derniers paramètres, jusqu'à quatre
appels \emph{par proposition}. Jusqu'à trois tours d'ajustement du champ peuvent proposer
de nouvelles initialisations. Ce n'est pas une modification du solveur en un appel de 120 itérations.

Le veto sur l'écart de gradient est supprimé pour ces essais. Une sortie retenue doit
avoir le statut convergé, une corrélation finale d'au moins 0,65, rester dans l'image
et conserver un déterminant de jacobienne positif sur tous les pixels échantillonnés
du support. \textbf{Ce sont des candidats convergés ; leur exactitude mécanique n'est pas
certifiée faute de vérité de déplacement associée.}

\newpage
\section*{1. Même région matérielle, supports et inconnues différents}
Pour un centre $\mathbf x$, la transformation locale est écrite
\[
 \boldsymbol\varphi(\mathbf x+\boldsymbol\xi)=\mathbf x+\mathbf d
 +(I+G)\boldsymbol\xi+
 \frac12\begin{pmatrix}
 H_{u,xx}\xi_x^2+2H_{u,xy}\xi_x\xi_y+H_{u,yy}\xi_y^2\\
 H_{v,xx}\xi_x^2+2H_{v,xy}\xi_x\xi_y+H_{v,yy}\xi_y^2
 \end{pmatrix}.
\]
L'affine possède six paramètres : deux déplacements et quatre dérivées premières.
Le quadratique ajoute six dérivées secondes. Il permet au gradient de varier dans
la fenêtre. Deux initialisations de la courbure sont comparées : zéro, ou dérivées
secondes analytiques du même champ RBF qui fournit le déplacement et le gradient.

\textbf{Subdivision réellement dense.} Chaque ancienne fenêtre non retenue de
$32\times32$ reçoit neuf centres fins aux décalages $\{-8,0,8\}^2$ pixels, avec un
support de $16\times16$. Les supports fins restent entièrement dans le support initial.
L'ancien centre est conservé exactement. Les quatre supports aux décalages $(\pm8,\pm8)$
pavent l'ancienne fenêtre ; les cinq autres ajoutent du recouvrement.

Il y a ainsi $327\times9=2\,943$ propositions fines, mais toujours seulement 327 anciens
centres à récupérer. Parmi les 2\,843 points fins retenus après continuation, 327
coïncident avec ces anciens centres et \textbf{2\,516 sont des centres supplémentaires}.
La superficie couverte est mesurée comme l'union des supports retenus, pas comme leur
somme avec recouvrements.

\textbf{Implémentation.} Les trois modèles s'exécutent avec le solveur Gauss--Newton
à composition inverse (IC-GN) existant, sur WebGPU. Le modèle quadratique à douze
paramètres était déjà disponible dans Live-DIC ; aucun code du solveur externe n'a
été modifié. Son API refuse le masque de région d'intérêt pour le quadratique :
le programme active donc explicitement les mêmes fenêtres à support plein que le
modèle affine. Il ne demande jamais de support partiel au solveur quadratique.

La sélection des nouvelles sorties n'applique plus le contrôle de gradient voisin,
ni la confirmation topologique de l'ancienne chaîne. L'écart avec les archives ne doit
donc pas être attribué au seul changement de support ou d'ordre polynomial.

\newpage
\section*{2. Le nombre de fenêtres et le budget de résolution ne disent pas la même chose}
\fig{comparison.pdf}{87mm}
Avec un appel par proposition, l'affine fine couvre 93,60\,\% de la superficie auparavant
en échec, contre 90,21\,\% pour l'affine 32. Elle ne récupère pourtant pas davantage
d'anciens centres : 294 contre 295. Les deux ensembles sont complémentaires.

La sélection déterministe « affine 32 si retenue, sinon affine 16 au même centre »
fournit \textbf{327/327 centres} : 295 avec support 32 et 32 avec support 16. Cet assemblage
utilise uniquement les décisions enregistrées, sans vérité analytique ni interpolation
des centres manquants. Il est exporté dans \texttt{mixed\_centres/}, avec la taille de
support et l'origine de chaque état. Il n'a pas subi une nouvelle résolution après assemblage.

La continuation retrouve aussi 327/327 centres en affine 32 seule, et en quadratique.
La grille fine atteint alors 98,38\,\% de couverture en superficie des régions traitées,
avec 100 propositions fines encore non retenues. Retrouver tous les anciens centres
ne signifie donc pas que toute la grille dense converge.

\newpage
\section*{3. Où les points supplémentaires sont-ils réellement calculés ?}
\fig{regions_191.pdf}{110mm}
Au niveau 191, cinq anciennes fenêtres restaient non retenues. Les rectangles orange
délimitent exactement ces cinq régions. Après continuation, les modèles affine et
quadratique 32 retrouvent leurs cinq centres ; la grille fine retient 41 points sur 45,
dont les cinq centres originaux.

Les points rouges sont conservés comme échecs. Ils ne sont pas remplacés par des
valeurs prédites présentées comme des observations. Les régions déjà retenues avant
cette expérience restent inchangées. Les figures ne montrent donc pas une nouvelle
grille dense sur toute la plaque : il s'agit d'une densification localisée aux échecs.

\newpage
\section*{4. Le quadratique améliore le recalage photométrique}
\fig{patch_65.pdf}{105mm}
La fenêtre 65 illustre la différence de modèle à support identique. L'affine 32 finit
par converger après continuation, mais le quadratique réduit davantage le résidu.
Le panneau affine 16 représente uniquement le carré central de 16 pixels : le gris
uniforme autour n'est ni une image corrigée ni une zone recalée.

Sur les 327 supports de 32 pixels, la corrélation médiane passe de \textbf{0,8463 en
affine à 0,9307 en quadratique}. Les minima sont respectivement 0,7607 et 0,8670.
Les deux initialisations du quadratique aboutissent à des valeurs presque identiques
après continuation. Le gain photométrique ne prouve pas, à lui seul, une meilleure
estimation de tous les paramètres mécaniques.

\newpage
\section*{5. Une fenêtre précédemment très mal initialisée}
\fig{patch_164.pdf}{105mm}
La fenêtre 164, qui conservait un mauvais résidu dans l'expérience précédente,
converge dans cette nouvelle chaîne : le champ est ajusté à mesure que de nouveaux
points sont résolus, le veto de gradient est absent, et la continuation laisse au
solveur davantage d'itérations. Le quadratique améliore encore la similitude du mouchetis.

Sur les 327 centres, les déplacements finaux affine et quadratique diffèrent de
0,703 pixel en médiane, et de 2,583 pixels au maximum. \textbf{Sans déplacements vrais,
ces écarts ne peuvent pas être interprétés comme des erreurs de l'un ou de l'autre.}
Ils justifient de conserver les modèles et supports dans la provenance de chaque graine.

\newpage
\section*{6. Stratégie retenue pour l'initialisation}
\textbf{1. Ne plus confondre incohérence avec le voisin et erreur de recalage.}
Dans cette expérience, le changement de gradient n'est plus un veto. Les statuts,
résidus photométriques, sorties hors image et inversions locales de la transformation
restent contrôlés. Les seuils de production n'ont pas été changés.

\textbf{2. Distinguer budget et représentation.} Une sortie non convergée après
30 itérations mérite une continuation bornée depuis son dernier état lorsqu'il est
fini. Cela a suffi à retrouver tous les anciens centres dans les 191 paires testées,
même en affine. Le modèle quadratique ne doit pas être jugé seulement au même nombre
d'itérations que l'affine ; ici, il converge plus difficilement dans la limite initiale.

\textbf{3. Densifier localement pour obtenir davantage de graines.} Les supports fins
avec pas réduit apportent de nouvelles observations. Leur association aux supports
32 récupère tous les anciens centres dès la campagne sans continuation. Il faut conserver
le support 16 dans les exports : une graine résolue sur 16 pixels n'est pas une mesure
validée sur une fenêtre de 32 pixels.

\textbf{4. Garder le quadratique pour les zones où l'affine laisse un résidu structuré.}
Sur des supports identiques, son amélioration photométrique est nette. Pour notre
objectif d'initialisation, il constitue une option complémentaire, pas une obligation
pour chaque fenêtre. La comparaison ne mesure pas encore la précision mécanique ni
le rendement d'une politique automatique de choix du modèle.

\textbf{Validation et limites.} Les 28 configurations pilotes relancées dans les campagnes
complètes reproduisent exactement leurs masques et paramètres. Treize tests ciblés
passent, notamment les dérivées secondes RBF par différences finies, l'ordre des douze
paramètres, le remappage quadratique et les positions des centres fins. Les deux campagnes
comptent chacune 764 configurations. Les résultats proviennent d'une seule plaque,
sans déplacements vrais associés. Aucune qualification de faux positifs n'est revendiquée.

Les scripts, états par fenêtre et par reprise, images, exports mixtes et empreintes
sont conservés dans \texttt{results/hudge\_adaptive\_support\_20261115/}.
Les essais infructueux de préparation sont aussi documentés : masque refusé par
l'API quadratique et demande de 100 itérations refusée par la limite de 30. La
continuation utilise l'API existante, sans modifier cette limite ni le solveur.
\end{document}
